 % Template:     Auxiliar LaTeX
% Documento:    Archivo principal
% Versión:      8.3.4 (07/02/2025)
% Codificación: UTF-8
%
% Autor: Pablo Pizarro R.
%        pablo@ppizarror.com
%
% Manual template: [https://latex.ppizarror.com/auxiliares]
% Licencia MIT:    [https://opensource.org/licenses/MIT]

% CREACIÓN DEL DOCUMENTO
\documentclass[
	spanish, % Idioma: spanish, english, etc.
	letterpaper, oneside
]{article}


\usepackage{amsmath}
\usepackage{breqn}

\newif\ifshowcomments
\showcommentstrue % Set to \showcommentsfalse to hide all comments

\usepackage{xcolor} % For colored text
\usepackage{marginnote} % Optional: for margin notes

% Define inline comment
\newcommand{\comment}[1]{%
  \ifshowcomments
    {\textcolor{red}{\textbf{ Solución:} }\\ #1 }%
  \fi
}
 

% INFORMACIÓN DEL DOCUMENTO
\def\documenttitle {Pauta Actividad 2 A2IC}
\def\documentsubject{Sesión RP de Modelamiento}

\def\documentauthor {Joaquín Ormazábal}
\def\coursename {Optimización}
\def\coursecode {EL4114}

\def\universityname {Universidad de Chile}
\def\universityfaculty {Facultad de Ciencias Físicas y Matemáticas}
\def\universitydepartment {Departamento de Ingeniería Eléctrica}
\def\universitydepartmentimage {departamentos/die}
\def\universitydepartmentimagecfg {height=1.75cm}
\def\universitylocation {Santiago de Chile}

% EQUIPO DOCENTE
\def\teachingstaff {
	\textbf{Profesor: Heraldo Rozas} \\
	Auxiliares: Diego Carez, Joaquín Ormazábal
}

% IMPORTACIÓN DEL TEMPLATE
\input{template}

% INICIO DE PÁGINAS
\begin{document}

% CONFIGURACIÓN DE PÁGINA Y ENCABEZADOS
\templatePagecfg

\fbox{\begin{minipage}{\textwidth}
\noindent\textbf{Instrucciones:} Formule los siguientes problemas de optimización, trabajando en grupos de 4 o 5 personas. Acuda al cuerpo docente en caso de consultas.
\end{minipage}}


% ======================= INICIO DEL DOCUMENTO ================
\begin{comment}
\newboxquestion{P1} \textbf{(Dificultad: Baja, Tiempo: 20 min)} Un observatorio astronómico debe planificar una noche de observación, decidiendo qué instrumentos científicos montar en el telescopio principal y qué observaciones realizar con cada uno, de modo de maximizar la utilidad científica.

El telescopio dispone de un conjunto $I$ de instrumentos posibles (por ejemplo, una cámara óptica, un espectrógrafo y una cámara infrarroja), los cuales deben montarse antes del inicio de la jornada. Por razones operativas y de espacio, sólo se pueden instalar a lo más $M$ instrumentos. Cada instrumento $i$ tiene un tiempo máximo de operación $\bar{t}_i$ (en horas), un costo fijo $f_i$ que se incurre si el instrumento se utiliza, y un costo variable $v_i$ por cada hora de operación. El observatorio dispone de un presupuesto total $B$ para cubrir los costos de operación de la noche.

Durante la noche pueden programarse observaciones pertenecientes al conjunto $O$, cada una asociada a distintos objetos astronómicos que pueden ser estudiados con algunos de los instrumentos disponibles. No todas las observaciones pueden realizarse con todos los instrumentos: para cada par $(o,i)$ se cuenta con un parámetro $a_{o,i}$ que indica si la observación $o$ puede efectuarse con el instrumento $i$. Además, si la observación $o$ se realiza con el instrumento $i$, requiere un tiempo de exposición $t_{o,i}$ (en horas). Cada observación aporta una utilidad científica $u_o$.

Plantee un problema de optimización que permita determinar qué instrumentos montar y qué observaciones realizar con cada uno, de modo de maximizar la utilidad científica total de la noche, respetando las limitaciones técnicas y económicas.
\end{comment}

\newboxquestion{P1} \textbf{(Tiempo: 45 min)} La empresa \textit{Maintenance} debe optimizar el calendario de reparaciones para una flota de $N_G$ turbinas de viento ubicadas en las costas de Inglaterra. 

Cada turbina $g \in \mathcal{G}=\{1,\ldots,N_G\}$ debe obligatoriamente recibir mantenimiento en una única fecha $t \in \mathcal{T}=\{1,\ldots,T_{\max}\}$ dentro de un horizonte de planificación. El tiempo de falla conocido de cada turbina es $\tau_g \in \mathcal{T}$. Si el mantenimiento se realiza antes de este tiempo de falla, la turbina continuará operando sin fallar durante todo el horizonte de planificación; sin embargo, si una turbina alcanza su tiempo de falla sin haber recibido mantenimiento, estará indisponible hasta que se le realice el mantenimiento correspondiente. Se asume que la reparación de cada turbina inicia y termina en tiempo tiempo $t$. Durante este periodo la turbina no está disponible.

Cuando una turbina $g$ está activa en el tiempo $t$, puede generar una potencia máxima conocida $P_{gt}[\mathrm{MWh}]$. No obstante, una turbina no está activa en tiempo $t$ si  ha fallado antes y aún no ha sido reparada, o se encuentra en mantenimiento en instante $t$. La energía generada por el parque eólico en cada instante $t$ es vendida a un precio unitario 
$\pi_t[\$/\mathrm{MWh}]$, el cual es conocido. La reparación de una turbina $g$ en el instante $t$ tiene un costo $C_{gt}$.

Las actividades de mantenimiento pueden ser calendarizadas solo cuando el equipo de mantenimiento está desplegado. El despliegue del equipo de mantenimiento en $t$ trae consigo un costo fijo adicional $A_t$. Notar que el equipo solo puede ser desplegado cuando las condiciones climáticas así lo permitan. En particular, se ha pronosticado que solo los días $\mathcal{K} \subset \mathcal{T}$ proveen condiciones seguras para el despliegue del equipo. Para un tiempo $t$, el equipo de mantenimiento tiene una capacidad máxima $M_t$. Por lo tanto, no se puede calendarizar la reparación de más de $M_t$ turbinas en tiempo $t$.


Existe una demanda de energía en cada instante de tiempo $t$, la cual conocida y se denota como $D_t$. Esta demanda constituye un requerimiento mínimo de generación, es decir el parque eólico en su conjunto debe producir al menos $D_t$.

Formule de forma ordenada un MILP que permita a la empresa \textit{Maintenance} coordinar las decisiones de operación y mantenimiento del parque eólico, de forma que se minimicen los costos netos. 


% PAUTA
\begin{comment}
{\color{MidnightBlue}
\vspace{5mm}

\textbf{Conjuntos e índices:}
\begin{align*}
I &: \text{Instrumentos posibles} \,(i \in I).\\
O &: \text{Observaciones posibles} \, (o \in O).
\end{align*}

\textbf{Parámetros:}
\begin{align*}
M &: \text{Cantidad máxima de instrumentos a instalar.}\\
\bar{t}_i &: \text{Tiempo máximo de operación del instrumento } i.\\
f_i &: \text{Costo fijo de utilizar el instrumento } i.\\
v_i &: \text{Costo variable por cada hora de operación del instrumento.}\\
B &: \text{Presupuesto de operación.}\\
a_{o,i} &: \text{Parámetro que indica si la observación } o \text{ necesita efectuarse con el instrumento } i.\\
t_{o,i} &: \text{Tiempo de exposición del instrumento } i \text{ necesario para realizar observación } o.\\
u_o &: \text{Utilidad científica de realizar la observación } o.
\end{align*}

\textbf{Variables:}
\begin{align*}
y_{o,i} &\in \{0,1\}: \text{ 1 si el instrumento } i \text{ se utiliza para la observación } o,\text{ 0 en caso contrario.}\\
x_i &\in \{0,1\}: \text{ 1 si el instrumento } i \text{ se instala antes del inicio de la jornada, 0 en caso contrario.}
\end{align*}

\textbf{Función objetivo:}
\begin{equation}
\max_{x_i, y_{o,i}} \quad \sum_{i\in I}\sum_{o\in O} u_o\,y_{o,i}
\end{equation}

\textbf{Restricciones:}

\begin{itemize}[label=\textcolor{MidnightBlue}{\textbullet}]
    \item Máxima cantidad de instrumentos instalados simultáneamente.
    \begin{equation}
        \sum_{i\in I} x_i \leq M
    \end{equation}
    \item Máximo tiempo de operación por instrumento.
    \begin{equation}
        \sum_{o\in O} t_{o,i}\,y_{o,i} \leq \bar{t}_i\,x_i \quad \forall i\in I
    \end{equation}
    \item El instrumento debe haberse instalado para poder ser utilizado.
    \begin{equation}
        y_{o,i} \leq x_i \quad \forall i\in I,\ \forall o\in O
    \end{equation}
    \item Cada observación se puede realizar a lo más una vez.
    \begin{equation}
        \sum_{i\in I} y_{o,i} \leq 1 \quad \forall o\in O
    \end{equation}
    \item Presupuesto disponible.
    \begin{equation}
        \sum_{i\in I}\left(f_i\,x_i+v_i\sum_{o\in O} y_{o,i}\right)\leq B
    \end{equation}
    \item Compatibilidad técnica observación-instrumento.
    \begin{equation}
        y_{o,i} \leq a_{o,i} \quad \forall i\in I,\ \forall o\in O
    \end{equation}
    \item Naturaleza de las variables.
    \begin{equation}
        y_{o,i},\,x_i \in \{0,1\} \quad \forall i\in I,\ \forall o\in O.
    \end{equation}
\end{itemize}
}
\end{comment}

\section*{Resolución P1}

Se nos solicita minimizar los costos netos de operación y mantenimiento del parque eólico,
que se definen por los (Costos totales de mantenimiento) - (Ingresos obtenidos por generación
de energía). Para esto, es conveniente definir las siguientes variables de decisión:

\begin{table}[htbp]
    \centering
    \small
    \caption{Variables de decisión}
    \label{tab:variables_decision}

    \begin{tabular}{p{0.15\textwidth} p{0.14\textwidth} p{0.55\textwidth}}
        \hline
        \textbf{Variable} & \textbf{Dominio} & \textbf{Descripción} \\
        \hline

        $x_{g,t}$ & $\{0,1\}$ &
        Indica si la turbina $g$ se encuentra en mantenimiento en el período $t$. \\

        $y_{g,t}$ & $\{0,1\}$ &
        Indica si la turbina $g$ se encuentra disponible en el período $t$. \\

        $z_t$ & $\{0,1\}$ &
        Indica si el equipo de mantenimiento se encuentra desplegado en el período $t$. \\

        $p_{g,t}$ & $\mathbb{R}^{+}$ &
        Indica la energía producida por la turbina $g$ en el período $t$. \\

        \hline
    \end{tabular}
\end{table}

Tras esto, podemos plantear la función objetivo del problema, la cual, tal como se mencionó
previamente, consiste en minimizar los costos netos, obteniéndose así que:

\begin{equation}
    FO:
    \min
    \sum_{g \in G}
    \sum_{t \in T}
    C_{g,t} \cdot x_{g,t}
    +
    \sum_{t \in T}
    A_t \cdot z_t
    -
    \sum_{g \in G}
    \sum_{t \in T}
    \pi_t \cdot p_{g,t}
    \tag{1}
\end{equation}

Donde el primer término representa los costos en los que se incurre por la reparación de cada turbina; el segundo término representa el costo asociado a desplegar el equipo de mantenimiento; mientras que el tercer término hace referencia a los ingresos que posee el parque eólico por la venta de energía. Luego, se definen las restricciones del problema, donde se consideran:

\begin{itemize}

    \item \textbf{Reparación única:} Fuerza a cada turbina a recibir mantenimiento en una
    única fecha.

    \begin{equation}
        \sum_{t \in T} x_{g,t} = 1,
        \quad \forall g \in G
    \end{equation}

    \item \textbf{Despliegue del equipo:} Fuerza a que en caso de que una turbina se encuentre
    siendo reparada en el período $t$, se despliegue el equipo de reparación.

    \begin{equation}
        x_{g,t} \leq z_t,
        \quad \forall g \in G,\; t \in T
    \end{equation}

    \item \textbf{Capacidad máxima:} Establece que no se puede calendarizar la reparación de
    más de $M_t$ turbinas en un tiempo $t$.

    \begin{equation}
        \sum_{g \in G} x_{g,t} \leq M_t,
        \quad \forall g \in G,\; t \in T
    \end{equation}

    \item \textbf{Mal clima:} Establece que el equipo solo se puede desplegar cuando las condiciones
    climáticas lo permiten, donde los días $K \subset T$ proveen las condiciones seguras.

    \begin{equation}
        z_t \leq 1 - \mathbf{1}_{\{t \notin \mathcal{K}\}},
        \quad \forall t \in T
    \end{equation}

    \item \textbf{Cumplir demanda:} Representa el requerimimiento mínimo de generación que
    debe tener el parque eólico.

    \begin{equation}
        \sum_{g \in G} p_{g,t} \geq D_t,
        \quad \forall t \in T
    \end{equation}

    \item \textbf{Límites de generación:} Establece que los límites técnicos de las turbinas.

    \begin{equation}
        0 \leq p_{g,t} \leq P_{g,t} \cdot y_{g,t},
        \quad \forall g \in G,\; t \in T
    \end{equation}

    \item \textbf{Disponibilidad:} Representa si las turbinas se encuentran activas o no. Una turbina
    puede estar no activa en el período $t$ si aún no supera el tiempo de falla $\tau_g$ pero se
    encuentra en mantención; o si ya ha superado el tiempo de falla, y aún no ha sido
    reparada.

    \begin{equation}
        y_{g,t} =
        \begin{cases}
            1 - x_{g,t}, & \text{si } t < \tau_g, \\[4pt]
            \displaystyle \sum_{r=1}^{t-1} x_{g,r}, & t \geq \tau_g.
        \end{cases}
    \end{equation}

\end{itemize}

%\newpage
\newboxquestion{P2} \textbf{(Tiempo: 45 min)} Una empresa de telecomunicaciones desea desplegar una red 5G en una ciudad. Para ello cuenta con un conjunto de sitios candidatos $J$ donde podrían instalarse estaciones base, y con un conjunto de zonas de demanda $K$, cada una con un número de usuarios $D_k$. El objetivo es determinar qué antenas instalar y en qué sitios, y cómo asignar la demanda de las distintas zonas a las antenas instaladas.

En cada sitio puede instalarse a lo más una antena, pudiendo esta ser de banda baja (tipo $L$), caracterizada por un mayor alcance pero menor capacidad, o de banda alta (tipo $H$), que tiene un alcance más reducido pero admite un volumen mayor de usuarios. Instalar una antena tipo $L$ en el sitio $j$ tiene un costo fijo $F_j^L$, mientras que instalar una antena tipo $H$ en el sitio $j$ tiene un costo fijo $F_j^H$. La empresa dispone de un presupuesto total $B$. Cada antena posee una capacidad nominal $C_j^L$ o $C_j^H$, dependiendo del tipo, que corresponde al número máximo de usuarios que puede atender en total.

Para cada zona $k$ y sitio $j$, los parámetros $a_{kj}^L$ y $a_{kj}^H$ indican si una antena tipo $L$ o $H$ instalada en $j$ podría atender usuarios de esa zona (por visibilidad y alcance). La demanda de una zona puede dividirse entre distintas antenas, y el modelo no exige que todas las zonas sean atendidas: se puede decidir no atender una zona. Una zona se considera \textit{atendida completamente} solo cuando la totalidad de su demanda $D_k$ es asignada a antenas instaladas con cobertura sobre ella. Únicamente en ese caso se obtiene una utilidad $u_k$ asociada a la zona. Es decir, no se genera utilidad parcial por atender solo una fracción de su demanda.

Además, la empresa desea que la red posea cierto grado de redundancia operativa. Para ello se considera que una zona queda ``bien cubierta'' únicamente cuando está atendida completamente y, además, al menos dos antenas instaladas, de cualquier tipo, tienen asignada una fracción mínima $\alpha$ de la demanda de dicha zona. La empresa exige que al menos $R$ zonas cumplan esta condición.

\begin{itemize}
    \item[a)] Formule un modelo de optimización lineal entero mixto que permita planificar la red maximizando las utilidades.
    \item[b)] Discuta qué ocurre con el valor óptimo del problema cuando se incrementa el requisito de zonas bien cubiertas, es decir, si $R$ aumenta.
\end{itemize}

% PAUTA
\section*{Resolución P2}
\textbf{Conjuntos e índices:}
\begin{align*}
J &: \text{Sitios candidatos } (j\in J).\\
K &: \text{Zonas de demanda } (k\in K).\\
T &= \{L,H\}: \text{Tipos de antenas.}
\end{align*}

\textbf{Parámetros:}
\begin{align*}
D_k &: \text{Demanda (número de usuarios) en la zona } k.\\
F_j^t &: \text{Costo fijo de instalación de la antena tipo } t \text{ en } j.\\
C_j^t &: \text{Capacidad nominal de la antena tipo } t \text{ en el sitio } j.\\
B &: \text{Presupuesto total.}\\
u_k &: \text{Utilidad por atender completamente la zona } k.\\
R &: \text{Mínimo número de zonas bien cubiertas.}\\
\alpha &: \text{Fracción mínima de la demanda } D_k \text{ requerida por antena para la condición de redundancia.}\\
a_{kj}^t &\in \{0,1\}: \text{Indicador de cobertura (1 si la antena } t \text{ en } j \text{ cubre la zona } k).\\
M &: \text{Constante Big-M.}
\end{align*}

\textbf{Variables de decisión:}
\begin{align*}
x_{jt} &\in \{0,1\}: \text{ 1 si se instala antena tipo } t \text{ en } j.\\
z_k &\in \{0,1\}: \text{ 1 si la zona } k \text{ es atendida completamente.}\\
w_k &\in \{0,1\}: \text{ 1 si la zona } k \text{ está bien cubierta.}\\
v_{kj} &\in \{0,1\}: \text{ 1 si } j \text{ cubre más de la fracción } \alpha \text{ de la demanda.}\\
d_{kj}^t &\geq 0: \text{ Cantidad de usuarios asignados a la antena tipo } t \text{ en el sitio } j.
\end{align*}

\textbf{Función objetivo:}
\begin{align}
\max \sum_{k\in K} u_k z_k
\end{align}

\textbf{Restricciones:}

\begin{itemize}
    \item Instalación única por sitio
    \begin{equation}
        \sum_{t\in T} x_{jt} \leq 1 \quad  \forall j\in J
    \end{equation}
    \item Cobertura
    \begin{equation}
        d_{kj}^t \leq M a_{kj}^t \quad \forall k\in K,\ \forall j\in J,\ \forall t\in T
    \end{equation}
    \item Capacidad por antena instalada
    \begin{equation}
        \sum_{k\in K} d_{kj}^t \leq C_j^t x_{jt} \quad \forall j\in J,\ \forall t\in T
    \end{equation}
    \item Restricción presupuestaria
    \begin{equation}
        \sum_{j\in J}\sum_{t\in T} F_j^t x_{jt} \leq B
    \end{equation}
    \item Satisfacción de demanda
    \begin{equation}
        \sum_{j\in J}\sum_{t\in T} d_{kj}^t \leq D_k \quad \forall k\in K
    \end{equation}
    \item Atención completa
    \begin{equation}
        \sum_{j\in J}\sum_{t\in T} d_{kj}^t \geq D_k z_k \quad \forall k\in K
    \end{equation}
    \item Condición de bien cubierta
    \begin{itemize}
        \item[a)] Condición de redundancia
        \begin{equation}
            \sum_{j\in J} v_{kj} \geq 2w_k \quad \forall k\in K
        \end{equation}
        \item[b)] Condición de atención completa
        \begin{equation}
            w_k \leq z_k \quad \forall k\in K
        \end{equation}
        \item[c)] Condición de fracción mínima
        \begin{equation}
            \sum_{t\in T} d_{kj}^t \geq \alpha D_k v_{kj} \quad  \forall k\in K,\ \forall j\in J
        \end{equation}
     \end{itemize}
    \item Requisito R
    \begin{equation}
        \sum_{k\in K} w_k \geq R
    \end{equation}
    \item Dominio de las variables
    \begin{equation}
        x_{jt},\,z_k,\,w_k,\,v_{kj} \in \{0,1\}, \quad d_{kj}^t \geq 0.
    \end{equation}
\end{itemize}


Por otro lado, al aumentar el valor de $R$ el problema se hace más restringido, con lo que el espacio factible del problema se reduce. Con esto, el óptimo será peor o igual que el problema original. Como es un problema de maximización, el nuevo valor óptimo será menor o igual al del problema original, pudiendo incluso resultar infactible.

\begin{comment}
\newboxquestion{P3} \textbf{(Dificultad: Alta, Tiempo: 40 min)} Una máquina puede realizar $m$ operaciones, indexadas de $1$ a $m$. Cada operación $j$ requiere una única herramienta $j$. La máquina puede alojar simultáneamente $B$ herramientas en su almacén de herramientas, donde $B < m$. Cargar o descargar la herramienta $j$ en el almacén de la máquina requiere $s_j$ unidades de tiempo de preparación. Sólo se puede cargar o descargar una herramienta cada vez.

Al comienzo del día, $n$ trabajos esperan a ser procesados por la máquina. Cada trabajo $i$ requiere varias operaciones. Sea $J_i$ el conjunto de operaciones requeridas por el trabajo $i$, y supongamos por simplicidad que para todo $i$, $|J_i|$ no es mayor que la capacidad del almacén $B$ de la máquina. Antes de que la máquina pueda empezar a procesar el trabajo $i$, todas las herramientas necesarias pertenecientes al conjunto $J_i$ deben
estar preparadas en la máquina. Si una herramienta $j \in J_i$ ya está cargada en la máquina, evitamos el tiempo de preparación de la herramienta $j$. Si la herramienta $j \in J_i$ no está cargada, debemos prepararla, posiblemente (si el almacén de herramientas está lleno) después de descargar una herramienta existente que el trabajo $i$ no necesita.

Una vez preparadas las herramientas, se procesan todas las operaciones $|J_i|$ del trabajo $i$. Hay que tener en cuenta que, debido a los requisitos comunes de las herramientas para los distintos trabajos y a la capacidad limitada del almacén, el tiempo de preparación necesario antes de cada secuencia de trabajo depende de cada trabajo. Formule un problema de programación entera para determinar la secuencia óptima de trabajos que minimice el tiempo total de preparación para completar todos los trabajos. Suponga que al principio del día, el almacén de herramientas está completamente vacío.

% PAUTA
{\color{MidnightBlue}
\textbf{Variables de decisión:}

\[
x_{ir} =
\begin{cases}
1, & \text{si el trabajo } i \text{ es el } r\text{-ésimo trabajo procesado},\\
0, & \text{en otro caso}.
\end{cases}
\]

\[
y_{ir} =
\begin{cases}
1, & \text{si la herramienta } j \text{ está en el almacén mientras el } r\text{-ésimo trabajo es procesado},\\
0, & \text{en otro caso}.
\end{cases}
\]

\textbf{Restricciones:}

\begin{itemize}
    \item Que el almacén esté vacío al inicio, se modela como
    \[
    y_{j0} = 0, \quad \forall j.
    \]

    \item Como todos los trabajos deben ser procesados, se requiere que
    \[
    \sum_{r=1}^{n} x_{ir} = 1, \quad \forall i.
    \]

    \item Como se debe procesar un trabajo a la vez, se debe cumplir que
    \[
    \sum_{i=1}^{n} x_{ir} = 1, \quad \forall r.
    \]

    \item Un trabajos $i$ se puede procesar sólo si las herramientas en $J_i$ están en el almacén, por lo tanto se debe incluir una restricción de forzamiento:
    \[
    x_{ir} \leq y_{ir}, \quad \forall j \in J_i, \ \forall r, i.
    \]

    \item El almacén tiene una capacidad máxima, así que
    \[
    \sum_{j=1}^{m} y_{ir} \leq B, \quad \forall r.
    \]
\end{itemize}

\textbf{Función objetivo:}

Los tiempos que se pueden gestionar son aquellos producidos por la carga y descarga de herramientas en el almacén, $s_j$, los cuales son utilizados cuando la herramienta requerida en la secuencia actual no estaba cargada en la secuencia anterior. Por lo tanto, la función objetivo es:

\[
\min \sum_{j=1}^{m} \sum_{r=1}^{n} s_j \left|y_{jr} - y_{j,r-1}\right|
\]

Esta función objetivo se puede reformular definiendo una nueva variable de decisión $z_{jr}$, quedando la función objetivo como:

\[
\min \sum_{j=1}^{m} \sum_{r=1}^{n} s_j z_{jr}
\]

Y agregando las restricciones

\[
z_{jr} \geq y_{jr} - y_{j,r-1}
\]

y

\[
z_{jr} \geq y_{j,r-1} - y_{jr}, \quad \forall j, r.
\]

\textbf{Modelo completo:}

\begin{align*}
\min \quad & \sum_{j=1}^{m} \sum_{r=1}^{n} s_j z_{jr} \\
\text{s.a.} \quad
& z_{jr} \geq y_{jr} - y_{j,r-1}, && \forall j, r \\
& z_{jr} \geq y_{j,r-1} - y_{jr}, && \forall j, r \\
& y_{j0} = 0, && \forall j \\
& \sum_{r=1}^{n} x_{ir} = 1, && \forall i \\
& \sum_{i=1}^{n} x_{ir} = 1, && \forall r \\
& x_{ir} \leq y_{ir}, && \forall j \in J_i, \ \forall r, i \\
& \sum_{j=1}^{m} y_{ir} \leq B, && \forall r \\
& x_{ir}, y_{jr}, z_{jr} \in \{0,1\}.
\end{align*}

}
\end{comment}

% FIN DEL DOCUMENTO
\end{document}